\documentclass[conference]{IEEEtran}
\IEEEoverridecommandlockouts

\usepackage{cite}
\usepackage{amsmath,amssymb,amsfonts}
\usepackage{algorithm}
\usepackage{algpseudocode}
\usepackage{graphicx}
\usepackage{textcomp}
\usepackage{xcolor}
\usepackage{hyperref}
\usepackage[utf8]{inputenc}
\usepackage[T1]{fontenc}
\usepackage{float}

\def\BibTeX{{\rm B\kern-.05em{\sc i\kern-.025em b}\kern-.08em
    T\kern-.1667em\lower.7ex\hbox{E}\kern-.125emX}}

\begin{document}

\title{Simulação de um Relógio Digital com Sistema de Alarme}

\author{\IEEEauthorblockN{Francisco Passos dos Santos Alves}
\IEEEauthorblockA{\textit{Department of Computer Science} \\
\textit{Federal University of Sergipe}\\
Aracaju, Sergipe, Brazil \\
francisco.alves@dcomp.ufs.br}

\and

\IEEEauthorblockN{João Vitor Vital Leão}
\IEEEauthorblockA{\textit{Department of Computer Science} \\
\textit{Federal University of Sergipe}\\
Aracaju, Sergipe, Brazil \\
joao.leao@dcomp.ufs.br}

\and

\IEEEauthorblockN{Lucca Pedreira Dultra}
\IEEEauthorblockA{\textit{Department of Computer Science} \\
\textit{Federal University of Sergipe}\\
Aracaju, Sergipe, Brazil \\
lucca.dultra@dcomp.ufs.br}

\and

\IEEEauthorblockN{Karlos Danyel Santos Gomes}
\IEEEauthorblockA{\textit{Department of Computer Science} \\
\textit{Federal University of Sergipe}\\
Aracaju, Sergipe, Brazil \\
karlos.gomes@dcomp.ufs.br}
}

\maketitle

\begin{abstract}
Este trabalho apresenta o desenvolvimento de um simulador computacional de relógio digital baseado em um projeto de hardware real criado por Wagner Rambo e disponibilizado em seu canal no YouTube, WR Kits. O simulador, denominado \textit{A Digital Clockwork Simulator}, foi implementado em linguagem C++ utilizando o paradigma de Programação Orientada a Objetos (POO), com o objetivo de reproduzir o comportamento lógico dos circuitos integrados presentes no projeto original, incluindo componentes como CD4017, CD4029, CD4511, CD4013, um gerador de frequência, entre outros. A motivação central do projeto é educacional: por meio da simulação em software, torna-se possível estudar o comportamento de circuitos digitais discretos, explorar conceitos de hardware de baixo nível e validar o funcionamento lógico do circuito sem a necessidade de implementação física. Como extensão do escopo original, o projeto incorpora um sistema de alarme desenvolvido como requisito adicional da disciplina de Sistemas Digitais. O projeto é distribuído sob a licença GNU General Public License (GPL) e está disponível publicamente no GitHub.
\end{abstract}

\begin{IEEEkeywords}
digital electronics, logic design, computer architecture, hardware modeling, simulation, educational project
\end{IEEEkeywords}


\section{Introdução}

\section{Fundamentação Teórica}

% \subsection{Eletrônica Digital}
% \subsection{Lógica Combinacional}
% \subsection{Lógica Sequencial}
% \subsection{Contadores e Divisores de Frequência}


\section{Circuito Original}

% \subsection{Descrição do Projeto de Referência}
% \subsection{Componentes Utilizados}
% \subsection{Funcionamento Geral do Circuito}


\section{Modelagem Computacional}

% \subsection{Objetivos da Simulação}
% \subsection{Representação dos Componentes}
% \subsection{Estratégia de Simulação}


\section{Arquitetura do Software}

% \subsection{Estrutura do Projeto}
% \subsection{Classes e Módulos}
% \subsection{Fluxo de Execução}


\section{Implementação dos Componentes}

% \subsection{CD4017}
% \subsection{CD4029}
% \subsection{CD4511}
% \subsection{Gerador de Frequência}
% \subsection{Sistema de Alarme}


\section{Testes e Validação}

% \subsection{Metodologia de Testes}
% \subsection{Testes Unitários}
% \subsection{Validação dos Resultados}


\section{Resultados}

% \subsection{Comportamento dos Contadores}
% \subsection{Comportamento dos Displays}
% \subsection{Funcionamento do Alarme}


\section{Limitações do Simulador}


\section{Trabalhos Futuros}


\section{Conclusão}


\section*{Agradecimentos}
Agradecemos ao Prof. Dr. Calebe Micael de Oliveira Conceição e ao Prof. Dr. Rodolfo Botto de Barros Garcia pelas orientações, ensinamentos e apoio durante o desenvolvimento deste trabalho. Em especial, o autor Francisco Passos dos Santos Alves agradece aos docentes por permitirem a utilização e continuidade da base do projeto desenvolvida em etapas anteriores da disciplina pelo referido autor, possibilitando sua expansão e aperfeiçoamento para a avaliação da terceira unidade de Sistemas Digitais. Agradecemos também à Universidade Federal de Sergipe (UFS) e ao Departamento de Computação pela infraestrutura disponibilizada.

\begin{thebibliography}{99}

\bibitem{ref1}
A. Lastname, B. Lastname, \textit{Book Title}, Xth ed. City: Publisher, Year.

\end{thebibliography}

\end{document}